\begin{thebibliography}{99}
\small
\addcontentsline{toc}{chapter}{Bibliography}

\bibitem{aft2011} A.~Atserias, J.~K. Fichte, M.~Thurley. Clause-learning algorithms with many restarts and bounded-width resolution. \emph{Journal of Artificial Intelligence Research} 40 (2011), 353--373.
\bibitem{bachmairganzinger1994} L.~Bachmair, H.~Ganzinger. Rewrite-based equational theorem proving with selection and simplification. \emph{Journal of Logic and Computation} 4(3) (1994), 217--247.
\bibitem{bdh1986} L.~Bachmair, N.~Dershowitz, J.~Hsiang. Orderings for equational proofs. \emph{Proc.\ 1st IEEE Symposium on Logic in Computer Science} (1986), 346--357.
\bibitem{birkhoff1935} G.~Birkhoff. On the structure of abstract algebras. \emph{Proceedings of the Cambridge Philosophical Society} 31 (1935), 433--454.
\bibitem{buckingham1914} E.~Buckingham. On physically similar systems; illustrations of the use of dimensional equations. \emph{Physical Review} 4(4) (1914), 345--376.
\bibitem{caccioppoli1932} R.~Caccioppoli. Sugli elementi uniti delle trasformazioni funzionali. \emph{Rendiconti del Seminario Matematico della Universit\`a di Padova} 3 (1932), 1--15.
\bibitem{church1936} A.~Church. An unsolvable problem of elementary number theory. \emph{American Journal of Mathematics} 58 (1936), 345--363.
\bibitem{cook1971} S.~A. Cook. The complexity of theorem-proving procedures. \emph{Proc.\ 3rd ACM Symposium on Theory of Computing} (1971), 151--158.
\bibitem{cookreckhow1979} S.~A. Cook, R.~A. Reckhow. The relative efficiency of propositional proof systems. \emph{Journal of Symbolic Logic} 44(1) (1979), 36--50.
\bibitem{coverthomas2006} T.~M. Cover, J.~A. Thomas. \emph{Elements of Information Theory}, 2nd ed. Wiley, 2006.
\bibitem{evans-pde} L.~C. Evans. \emph{Partial Differential Equations}, 2nd ed. American Mathematical Society, 2010.
\bibitem{farkas1902} J.~Farkas. Theorie der einfachen Ungleichungen. \emph{Journal f\"ur die reine und angewandte Mathematik} 124 (1902), 1--27.
\bibitem{haken1985} A.~Haken. The intractability of resolution. \emph{Theoretical Computer Science} 39 (1985), 297--308.
\bibitem{hadamard1906} J.~Hadamard. Sur les transformations ponctuelles. \emph{Bulletin de la Soci\'et\'e Math\'ematique de France} 34 (1906), 71--84.
\bibitem{herbrand1930} J.~Herbrand. \emph{Recherches sur la th\'eorie de la d\'emonstration}. Thesis, Universit\'e de Paris, 1930.
\bibitem{huet1981} G.~Huet. A complete proof of correctness of the Knuth--Bendix completion algorithm. \emph{Journal of Computer and System Sciences} 23 (1981), 11--21.
\bibitem{knuthbendix1970} D.~E. Knuth, P.~B. Bendix. Simple word problems in universal algebras. In: J.~Leech (ed.), \emph{Computational Problems in Abstract Algebra}, Pergamon, 1970, 263--297.
\bibitem{krantzparks2013} S.~G. Krantz, H.~R. Parks. \emph{The Implicit Function Theorem: History, Theory, and Applications}. Birkh\"auser, 2013.
\bibitem{eliashbergmishachev2002} Y.~Eliashberg, N.~Mishachev. \emph{Introduction to the $h$-Principle}. American Mathematical Society, 2002.
\bibitem{kirszbraun1934} M.~D. Kirszbraun. \"{U}ber die zusammenziehende und Lipschitzsche Transformationen. \emph{Fundamenta Mathematicae} 22 (1934), 77--108.
\bibitem{kuiper1955} N.~H. Kuiper. On $C^1$-isometric imbeddings I, II. \emph{Indagationes Mathematicae} 17 (1955), 545--556, 683--689.
\bibitem{nash1954} J.~Nash. $C^1$ isometric imbeddings. \emph{Annals of Mathematics} 60 (1954), 383--396.
\bibitem{tarski1951} A.~Tarski. \emph{A Decision Method for Elementary Algebra and Geometry}, 2nd ed. University of California Press, 1951.
\bibitem{nelsonoppen1979} G.~Nelson, D.~C. Oppen. Simplification by cooperating decision procedures. \emph{ACM Transactions on Programming Languages and Systems} 1(2) (1979), 245--257.
\bibitem{not2006} R.~Nieuwenhuis, A.~Oliveras, C.~Tinelli. Solving SAT and SAT modulo theories: from an abstract Davis--Putnam--Logemann--Loveland procedure to DPLL($T$). \emph{Journal of the ACM} 53(6) (2006), 937--977.
\bibitem{pd2011} K.~Pipatsrisawat, A.~Darwiche. On the power of clause-learning SAT solvers as resolution engines. \emph{Artificial Intelligence} 175(2) (2011), 512--525.
\bibitem{robinson1965} J.~A. Robinson. A machine-oriented logic based on the resolution principle. \emph{Journal of the ACM} 12(1) (1965), 23--41.
\bibitem{rockafellar1970} R.~T. Rockafellar. \emph{Convex Analysis}. Princeton University Press, 1970.
\bibitem{rogers1967} H.~Rogers, Jr. \emph{Theory of Recursive Functions and Effective Computability}. McGraw-Hill, 1967.
\bibitem{rudin1976} W.~Rudin. \emph{Principles of Mathematical Analysis}, 3rd ed. McGraw-Hill, 1976.
\bibitem{schrijver1986} A.~Schrijver. \emph{Theory of Linear and Integer Programming}. Wiley, 1986.
\bibitem{tinelliharandi1996} C.~Tinelli, M.~T. Harandi. A new correctness proof of the Nelson--Oppen combination procedure. In: \emph{Frontiers of Combining Systems} (FroCoS 1996), Kluwer, 1996, 103--119.
\bibitem{tseitin1968} G.~S. Tseitin. On the complexity of derivation in propositional calculus. In: A.~O. Slisenko (ed.), \emph{Studies in Constructive Mathematics and Mathematical Logic, Part II}, 1968, 115--125.
\bibitem{turing1937} A.~M. Turing. On computable numbers, with an application to the Entscheidungsproblem. \emph{Proceedings of the London Mathematical Society} (2) 42 (1937), 230--265.

\medskip\noindent\textbf{Manuscripts of the author.}
\bibitem{flyxion-binding} Flyxion. \emph{The Algebra of Binding: Distinction, Interaction, and Admissible Collapse from Elementary Identities to General Systems}. Manuscript, 2026.
\bibitem{flyxion-disposition} Flyxion. \emph{The Calculus of Disposition}. Manuscript in preparation, 2026.
\bibitem{flyxion-foundations} Flyxion. \emph{Rebuilding the Foundations of Logic: Sets, Type Theory, and Reproducibility}. Manuscript, 2026.
\bibitem{flyxion-grounding} Flyxion. \emph{Grounding Before Collapse: Formulation Gaps, Verifier Gates, and the Limits of Self-Contained Judgment}. Manuscript, 2026.
\bibitem{flyxion-mutation} Flyxion. Verification inheritance and the Lean mutation harness. Manuscripts and tooling, 2026. (Descriptive title.)
\bibitem{flyxion-mips} Flyxion. \emph{Capacity, Dissipation, and the Onset of Motility-Induced Phase Separation}. Manuscript, 2026.
\end{thebibliography}
